\documentclass[10pt, legalpaper]{article}

\usepackage[a4paper, top=0.3in, bottom=0.3in, left=0.4in, right=0.4in]{geometry}

\usepackage[T1]{fontenc}
\usepackage{lmodern}
\usepackage[utf8]{inputenc}
\usepackage{enumitem}
\usepackage{titlesec}
\usepackage{hyperref}
\usepackage[expansion=false]{microtype}
\usepackage{xcolor}
\usepackage{array}
\usepackage{tabularx}

% --- COLORS ---
\definecolor{accent}{RGB}{0, 0, 0}
\definecolor{linkblue}{RGB}{0, 70, 127}
\definecolor{darkgray}{RGB}{35, 35, 35}
\definecolor{medgray}{RGB}{70, 70, 70}

% --- HYPERREF ---
\hypersetup{
    colorlinks=true,
    urlcolor=linkblue,
    linkcolor=linkblue,
    pdftitle={MD Humayun - Resume},
    pdfauthor={MD Humayun},
    pdfsubject={Embedded and Systems Engineer}
}

% --- SECTION FORMATTING ---
\titleformat{\section}
  {\normalsize\bfseries\color{accent}}
  {}
  {0em}
  {\MakeUppercase}
  [\vspace{-7pt}\color{accent}\rule{\linewidth}{0.8pt}\vspace{-3pt}]

\titlespacing{\section}{0pt}{4pt}{2.5pt}

% --- LIST SETTINGS ---
\setlist[itemize]{
    leftmargin=1.3em,
    topsep=0pt,
    itemsep=1.5pt,
    parsep=0pt,
    partopsep=0pt,
    label=$\bullet$
}

% --- CUSTOM COMMANDS ---
\newcommand{\eduentry}[4]{%
    \noindent\textbf{#1} \hfill \textit{\color{medgray}#2}\\[1pt]
    \noindent\textit{\color{darkgray}#3} \hfill \textbf{#4}\\[2pt]
}

\newcommand{\expentry}[3]{%
    \noindent\textbf{#1} \hfill {\small #3}\\[1pt]
    \noindent\textit{#2}\\[2pt]
}

\newcommand{\projentry}[4]{%
    \noindent\textbf{#1} \hfill \textit{\small\color{medgray}#2}\\[1pt]
    \noindent{\small\textit{\color{darkgray}#3}} \hfill \href{#4}{\small [GitHub]}\\[2pt]
}

% ============================================================
\begin{document}
\raggedbottom
\tolerance=1200
\emergencystretch=2em
\hbadness=10000
\pagestyle{empty}

% ============================================================
%  HEADER
% ============================================================
\begin{center}
{\LARGE\bfseries MD HUMAYUN}\\[2pt]
{\small
+91-7294963049 \textbar\
\href{mailto:md.humayun7777777@gmail.com}{md.humayun7777777@gmail.com} \textbar\
\href{https://www.linkedin.com/in/md-humayun-82051521a/}{LinkedIn} \textbar\
\href{https://github.com/mdhumayun7}{GitHub} \textbar\
\href{https://mdhumayun7.github.io/MD-HUMAYUN-PORTFOLIO/}{Portfolio} \textbar\
\href{https://leetcode.com/u/mdhumayun7/}{LeetCode} \textbar\
Surat, India
}\\[3pt]
{\small\textbf{Embedded \& Systems Engineer \textbar\ ESP32 \textbar\ FreeRTOS \textbar\ C++ \textbar\ MQTT \textbar\ GATE CS 93.8 Percentile}}
\end{center}

\vspace{-2pt}

% ============================================================
%  SUMMARY
% ============================================================
\section{Summary}

{\small
M.Tech CS student at SVNIT Surat (GATE \textbf{93.8} percentile) building systems where the part that decides never waits on the part that reports: ESP32 firmware with safety logic and networking on separate cores, tested on the host in CI. Also built a multi-agent drone swarm simulator under a MeitY-sponsored project, with a paper under review at IEEE GSCON 2027, and an SDN DDoS mitigator measured under live attack traffic.
}

% ============================================================
%  EDUCATION
% ============================================================
\section{Education}

\eduentry
{Sardar Vallabhbhai National Institute of Technology (SVNIT Surat)}
{2025 -- Present}
{M.Tech in Computer Science (Information Security \& Privacy) \textbar\ \textbf{CPI: 8.70/10}}
{Surat, Gujarat}

\vspace{-2mm}
\eduentry
{Bihar Engineering University}
{2020 -- 2024}
{B.Tech in Computer Science Engineering \textbar\ \textbf{CGPA: 8.51/10}}
{Patna, Bihar}

% ============================================================
%  EXPERIENCE
% ============================================================
\vspace{-2mm}
\section{Professional Experience}

\expentry
{Research Intern --- Hybrid AI-Based Autonomous Drone Swarm Navigation}
{MeitY (Ministry of Electronics \& IT), Govt. of India --- Sponsored Project, SVNIT Surat (Guide: \href{https://svnit.ac.in/maz/}{Prof. Mukesh A. Zaveri})}
{May 2026 -- Jul 2026}
\vspace{-2pt}
\noindent{\small\textit{\color{darkgray}Python \textbar\ Blender \textbar\ LLMs (Groq) \textbar\ LoRA \textbar\ PPO \textbar\ Gymnasium}} \hfill {\small\textit{\color{medgray}Team project (2)} \textbar\ \href{https://drive.google.com/file/d/1zHO33qPGtrWtDfhehBWT-QbK_ZDj5pvY/view?usp=sharing}{[Demo Video]}}\\[2pt]
\vspace{-4mm}
\begin{itemize}
\item Designed \textbf{BattleWorld}, a multi-agent simulator wrapped as a Gymnasium environment with a uniform \textbf{12-dimensional} state interface, so every commander --- classical or learned --- consumed identical input and the benchmark measured policy rather than input.
\item Built the scenario pipeline: procedural urban scenes in Blender auto-labelled through an LLM annotation stage, making scenario count a parameter instead of a bottleneck.
\item Fine-tuned an LLM with LoRA on the generated tactical labels, reaching \textbf{91.2\%} mean token accuracy at \textbf{0.463} training loss.
\item Benchmarked \textbf{seven} commander architectures over \textbf{28} held-out scenarios; hybrid LLM+PPO reached \textbf{67.9\%} action accuracy against \textbf{60.7\%} for the best classical baseline.
\end{itemize}

\vspace{1pt}

\expentry
{Salesforce Developer Intern}
{DST (Digital Skills Training) --- Remote}
{Dec 2023 -- Jan 2024}
\vspace{-4mm}
\begin{itemize}
\item Developed Apex triggers and automated workflows in OOP-based backend logic, cutting manual data-entry overhead by \textbf{35\%}.
\item Customised KPI dashboards for \textbf{5+} cross-functional teams, replacing manually compiled weekly reports.
\end{itemize}

% ============================================================
%  PUBLICATIONS
% ============================================================
\section{Publications}

{\small
\noindent\textbf{Hybrid AI-Based Autonomous Drone Swarm Navigation} \textbar\ \textit{\color{darkgray}Guided by Prof. Mukesh A. Zaveri, SVNIT Surat} \hfill \textit{\color{medgray}Under Review, IEEE GSCON 2027 (\#71435)}
}

% ============================================================
%  PROJECTS
% ============================================================
\section{Projects}

\projentry
  {Smart Home Surveillance --- ESP32 Alarm Firmware \textnormal{\small(v2.1)}}
  {2024 -- 2026}
  {C++ \textbar\ ESP32 \textbar\ FreeRTOS \textbar\ MQTT \textbar\ Wokwi \textbar\ GitHub Actions}
  {https://github.com/mdhumayun7/Smart-Home-Automation}
\vspace{-4mm}
\begin{itemize}
\item Split the firmware across the ESP32's two cores: the safety loop (sensor scan, arm/disarm state machine with \textbf{10 s} exit delay, \textbf{10 s} entry delay and \textbf{30 s} alarm, buzzer, LEDs, OLED) on core 1 with no blocking \texttt{delay()}, and Wi-Fi, NTP and MQTT as a FreeRTOS task on core 0, joined only by two queues, so a stalled network call cannot delay the siren.
\item Reported state over MQTT with Last Will availability, retained state, sequence-numbered JSON events and a \textbf{16-deep} offline queue delivered in order after reconnect; remote DISARM is compiled out unless TLS and broker credentials are configured.
\item Persisted the armed state in NVS across power cuts, supervised the loop with the task watchdog, debounced door and button inputs at \textbf{40 ms} and filtered the PIR with a \textbf{150 ms} confirmation window.
\item Tested the state machine on the host in CI: \textbf{14} scenarios plus a TLS build check at three \texttt{millis()} start values, including 32-bit wrap-around, on arduino-esp32 2.x and 3.x (\textbf{84} runs, \texttt{-Wall -Wextra -Werror}).
\end{itemize}

\vspace{1pt}

\projentry
  {IPS-ML / ShieldNet --- Intrusion Prevention with Explainable AI \textnormal{\small(M.Tech Project)}}
  {Aug 2025 -- Present}
  {Python \textbar\ scikit-learn \textbar\ XGBoost \textbar\ SHAP \textbar\ CICIDS2017 \textbar\ pytest}
  {https://github.com/mdhumayun7/ips}
\vspace{-4mm}
\begin{itemize}
\item Trained six detectors and a soft-voting ensemble on \textbf{2,497,980} audited CICIDS2017 flows, evaluated on an identical held-out set of \textbf{749,394}; ensemble reached \textbf{0.8814} macro-F1 at a \textbf{0.00135} false alarm rate and \textbf{0.432 ms} median inference.
\item Replaced argmax with a confidence-aware ALLOW / UNCERTAIN / BLOCK decision, choosing the operating point by sweeping the threshold against a stated \textbf{1\%} false-alarm budget; routing \textbf{1.62\%} of traffic to review cut missed attacks to \textbf{0.0149\%} and benign disruption to \textbf{0.0486\%}.
\item Computed SHAP attributions inside the enforcement path and wrote them to an append-only audit log beside each action, so every block is explainable after the fact.
\item Improved rare-class detection over the published Random Forest baseline without synthetic oversampling: SQL injection F1 \textbf{0.11 $\rightarrow$ 0.50}, Heartbleed \textbf{0.77 $\rightarrow$ 1.00}; covered by \textbf{29} unit tests with all thirteen plan deviations logged.
\end{itemize}

\vspace{1pt}

\projentry
  {DDoS Detection \& Mitigation --- SDN + Machine Learning}
  {Aug 2025 -- Jan 2026}
  {Python \textbar\ Mininet \textbar\ POX \textbar\ OpenFlow \textbar\ Random Forest \textbar\ Scapy}
  {https://github.com/mdhumayun7/ddos-5g-sdn-project}
\vspace{-4mm}
\begin{itemize}
\item Implemented detection inside a POX SDN controller tracking per-host packet rate, installing OpenFlow \texttt{DROP} rules in real time so mitigation is a control-plane action rather than an external appliance.
\item Trained a Random Forest on SDN traffic features to replace the fixed-threshold baseline, which misfires on legitimate traffic spikes.
\item Measured impact with \texttt{iperf} and \texttt{ping} rather than asserting it: throughput fell \textbf{54.3 $\rightarrow$ 31 Gbps} and average latency rose over \textbf{10x} under an ICMP flood.
\end{itemize}

\vspace{1pt}

\projentry
  {Multi-Source Job Intelligence Platform}
  {Apr 2026 -- Present}
  {Python \textbar\ Playwright \textbar\ SQLite \textbar\ GitHub Actions \textbar\ Streamlit \textbar\ Groq}
  {https://github.com/mdhumayun7/job-agent}
\vspace{-4mm}
\begin{itemize}
\item Built a headless scraping pipeline across \textbf{11} job boards and \textbf{9} government portals, persisting \textbf{2,158} listings to SQLite with idempotent upserts and per-scraper retry isolation, so one broken source degrades that source instead of failing the run.
\item Applied fuzzy deduplication via \texttt{SequenceMatcher} at a tuned \textbf{0.88} threshold with weighted skill scoring, reducing \textbf{793} raw records to \textbf{357} ranked matches --- a \textbf{55\%} noise cut.
\item Made LLM enrichment optional rather than required: a Groq-hosted model re-scores listings when configured, with a rule-based scorer as fallback, so the pipeline never depends on a paid service to run.
\item Orchestrated an \textbf{8-stage} pipeline on GitHub Actions cron, fanning out to Excel workbooks, HTML digests, email alerts and a Streamlit dashboard with no manual trigger.
\end{itemize}

\vspace{1pt}

% ============================================================
%  TECHNICAL SKILLS
% ============================================================
\section{Technical Skills}

{\small
\noindent\textbf{Languages:} C++ (primary), Python, MySQL, Bash \\
\noindent\textbf{CS Core:} Data Structures \& Algorithms, System Design, OOP \& Design Patterns, Operating Systems, DBMS, Computer Networks, Cryptography \\
\noindent\textbf{ML / AI:} TensorFlow, Keras, Scikit-learn, XGBoost, SHAP, LoRA, PPO, CNN, LSTM, VAE, Reinforcement Learning \\
\noindent\textbf{LLM \& Retrieval:} LangChain, Pinecone, Ollama, Groq, RAG, Gymnasium \\
\noindent\textbf{Embedded \& IoT:} ESP32, FreeRTOS, MQTT, Arduino, Raspberry Pi, I2C and GPIO, Wokwi simulation, host-side unit testing \\
\noindent\textbf{Web / Tools:} Flask (REST API), Docker, Git \& GitHub, GitHub Actions, AWS (EC2, S3), Linux, SDN, OpenFlow, Mininet
}

% ============================================================
%  CERTIFICATIONS & ACHIEVEMENTS
% ============================================================
\section{Certifications \& Achievements}
{\small
\noindent\textbf{GATE CS 2025:} 93.8 Percentile (National Eligibility for M.Tech Admissions) \\
\noindent\textbf{NPTEL:} Cloud Computing \& Ethical Hacking --- \textbf{All India Rank 33} \\
\noindent\textbf{Summer Research Internship Certificate:} SVNIT Surat, MeitY-Sponsored Project, issued by Prof. M. A. Zaveri \\
\noindent\textbf{IIT Bombay:} Python for Data Science, Computer Fundamentals \textbar\ \textbf{IIT Roorkee:} IoT \& Embedded Systems \\
\noindent\textbf{Community:} Founded NITIAN RAHI --- engineering education community of \textbf{52,000+} across YouTube, Facebook and Instagram
}

\end{document}
